%==============================================================================
% Section 3: Methodology
% Verbatim from Nckh-official_2.docx (P29-P55).
% The source labels the subsections "A."--"E."; the labels are dropped because
% llncs.cls numbers them automatically as 3.1--3.5 (the same numbering the
% source's own cross-references use).
% The bold run-in headings of the source ("Patch Encoding:", ...) are rendered
% with \subsubsection, the LNCS third-level run-in heading.
%==============================================================================
\section{Methodology}
\label{sec:method}

% ---- 3.1 --------------------------------------------------------------------
\subsection{Dataset and Evaluation Protocol}
\label{subsec:dataset}

The present work utilises the lithium-ion battery lifetime dataset from NASA
Ames Prognostics Center of Excellence~\cite{saha2007}. The dataset consists of
18650-type cells with a nominal capacity of 2.0\,Ah cycled until the capacity
falls below a safety threshold. Each discharge cycle has multivariate
time-series data (voltage, current, surface temperature) and the end-of-cycle
measured capacity. State of Health (SOH) is defined as current capacity over
nominal capacity (2.0\,Ah) in percentage form.

Quality filtering resulted in the retention of 26 of the original 34 cells.
Excluded cells were B0036 (excessive noise), B0049--B0052 (incomplete
sequences) and B0038--B0040 (reserved). We further removed cycles with missing
capacity values or with too short length. This work uses a strict cross-battery
split protocol, which differs from previous work that usually splits data by
time step within the same cell. Specifically:
% NOTE: the source reads "... within the same cell. Namely: Specifically:" --
% a duplicated lead-in; only "Specifically:" is kept here.
\begin{itemize}
  \setlength{\itemsep}{0pt}\setlength{\parskip}{0pt}
  \item \textbf{Training set:} 23 cells,
  \item \textbf{Validation set:} 2 cells (B0046, B0047),
  \item \textbf{Testing set:} 1 cell (B0048), completely isolated from
        training.
\end{itemize}

All validation and test cells were operated at \degC{4}. To support learning in
this temperature regime, the training set included several \degC{4} cells
(B0041, B0045, B0053--B0056). Preprocessing and splitting used a fixed random
seed of 42 for reproducibility. For additional statistical robustness, a
leave-one-battery-out (LOBO) protocol was applied on the window-30 variant, with
scalers refitted in each fold to avoid leakage. Results are presented in
Section~\ref{subsec:lobo}.

% ---- 3.2 --------------------------------------------------------------------
\subsection{Data Preprocessing and Feature Extraction}
\label{subsec:preprocessing}

Each discharge cycle provides four raw channels (voltage, current, surface
temperature, time stamp) plus two derived channels: the Incremental Capacity
($\mathrm{d}Q/\mathrm{d}V$) curve, which captures aging through shifts in peak
amplitude and position~\cite{weng2016}, and discharge progress, the normalized
accumulated discharged capacity (range $[0, 1]$), providing an explicit
positional reference within the cycle. As the battery ages, degradation effects
appear across the entire discharge profile, including shortened duration,
shifted temperature, and systematic decline of the IC peak near 3.5\,V, showing
that aging signatures are distributed throughout the cycle rather than confined
to narrow windows and motivating the use of the full discharge profile
($L = 4096$) as model input.

All six channels are normalized to $[0, 1]$ using \texttt{MinMaxScaler}, fitted
exclusively on the training set to prevent data leakage. Full-cycle sequences
are resampled to $L = 4096$ steps, while the compact configuration uses
non-overlapping 30-step sliding windows for real-time inference
(Section~\ref{subsec:training}), employing a dedicated set of six channels (four
raw measurements, normalized cycle index, and estimated SOC).

Each cycle is additionally summarized into a 57-dimensional
spectral-statistical feature vector (10 frequency-domain and 9 time-domain
descriptors per channel; full list in the project repository), computed over the
full cycle for optimal FFT resolution, shared across sliding windows within the
same cycle, and $z$-score normalized on the training distribution.

% ---- 3.3 --------------------------------------------------------------------
\subsection{Architecture of MambaSOHPredictor}
\label{subsec:architecture}

\subsubsection{Core Mamba Processing in Pure PyTorch.}
The computational core is based on the Selective State Space Model
(Mamba)~\cite{gu2023} implemented in pure native PyTorch without any special
CUDA kernels, allowing cross-platform operation (including Windows) and direct
CPU inference.

% Figure 1 of the source manuscript. Caption below the illustration, as
% required by the Springer LNCS style.
\begin{figure}[H]
  \centering
  \includegraphics[width=0.95\textwidth]{architecture.png}
  \caption{Overall architecture of the proposed framework}
  \label{fig:architecture}
\end{figure}

Each MambaBlock consists of an expansion projection ($E = 2$), a depthwise
causal convolution (kernel $= 4$), and a selective SSM scan whose state
parameters are dynamically adapted from the input~\cite{gu2023}, which
selectively keeps or discards information at each step, suiting battery signals
where most of the cycle is stationary and ageing features are concentrated in
transitional phases. This leads to linear $O(L)$ complexity for the SSM block
compared to the quadratic $O(L^{2})$ cost of self-attention at $L = 4096$.

\subsubsection{Patch Encoding.}
The input sequence is divided into patches of 16 steps (stride 8, 50\% overlap)
and projected to around 511 tokens, similar to PatchTST~\cite{nie2023}. Each
token is augmented with four local degradation statistics (RMS, peak-to-peak,
standard deviation, kurtosis) using the PatchDegradationEncoder, providing
dynamic context on top of the global spectral features.

\subsubsection{Token Aggregation (Attention Pooling).}
After two MambaBlocks and LayerNorm, we aggregate the 511 tokens with
single-head attention pooling to a single cycle-level representation. This
learns a weight distribution for the most informative dynamic phases of the
discharge cycle.

\subsubsection{Representation Modulation via FiLM.}
The 57-dimensional spectral vector modulates the pooled representation via
FiLM~\cite{perez2018}: a two-layer MLP produces scale/shift parameters
$(\gamma, \beta)$ to transform $h$ as
$h \leftarrow (\sigma(\gamma) + 0.5)\odot h + \beta$, which allows the spectral
summary to modulate features directly rather than only at the output
(Section~\ref{sec:discussion}).
% NOTE: the source points to "Section 5.2" here, but Section 5 (Discussion) has
% no subsections. Retargeted to Section 5 as a whole. Please confirm.

\subsubsection{Final Regression Head.}
The aggregated representation is passed through a two-layer head
($64 \rightarrow 32 \rightarrow 1$, GELU, 0.3 dropout) to output the continuous
SOH value.

% ---- 3.4 --------------------------------------------------------------------
\subsection{Training Setup}
\label{subsec:training}

The long-sequence configuration is optimized with AdamW (lr $5\times10^{-4}$,
weight decay $3\times10^{-4}$) against a weighted SmoothL1 loss
($\beta = 0.02$, clipping gradients for residuals above 2\% SOH), with near-EOL
samples ($\mathrm{SOH} < 80\%$) upweighted up to $2\times$, the weight decaying
exponentially with distance from the threshold. Training uses a short warm-up
(\texttt{CosineAnnealingWarmRestarts}, $T_{0} = 3$) followed by a final stage
($T_{0} = 80$, $T_{\mathrm{mult}} = 2$) whose periodic restarts escape the
plateau seen under fixed decay. The compact window-30 configuration is trained
separately with Adam (lr $5\times10^{-4}$, weight decay $10^{-5}$), batch size
32, up to 100 epochs (patience $= 15$).

% ---- 3.5 --------------------------------------------------------------------
\subsection{Uncertainty Quantification and Anomaly Detection}
\label{subsec:uncertainty}

\subsubsection{Uncertainty Quantification.}
Predictive uncertainty is estimated via Monte Carlo Dropout: 10 stochastic
forward passes (dropout active at inference, batched for latency) yield a mean
SOH estimate and standard deviation as the confidence band, with negligible
overhead given the compact model's 79{,}467 parameters.

\subsubsection{Anomaly Detection and Risk Management.}
An Isolation Forest module~\cite{liu2008} (100 estimators, contamination
$= 0.1$, seed $= 42$), fitted on the 57-dimensional spectral feature space,
classifies states via \texttt{decision\_function} into Normal, Warning
($\le -0.1$), and Anomaly ($\le -0.3$), orthogonal to four SOH-derived levels:
Healthy ($\ge 90\%$), Degrading ($\ge 85\%$), Maintenance Required
($\ge 80\%$), End of Life ($< 80\%$), per NASA guidelines.

These two axes are combined with hardware safety thresholds via a rule-based
risk matrix that assigns maintenance priority tiers, preserving a decoupled
design where deep learning handles long-term trends, Isolation Forest detects
sensor disturbances, and physical constraints act as non-overridable safeguards.
